%=======================================================================
%  REINVENT4 / Boltz-2 Closed-Loop Generative Pipeline
%  MANUSCRIPT TEMPLATE  --  Summer/Fall 2026
%
%  HOW TO USE THIS FILE
%  --------------------
%  1. Compile with: pdflatex -> bibtex -> pdflatex -> pdflatex
%     (or just hit "Recompile" on Overleaf).
%  2. Coloured boxes are INSTRUCTIONS, not manuscript text. They tell you
%     what to write, what to look up, and what figures to make.
%       * BLUE  "Guidance"  = what this section must accomplish
%       * ORANGE "Your turn" = you personally write this part
%       * GREEN "Figure"     = recommended figure, with the script that makes it
%  3. When you are ready to produce a clean draft for a reader, change
%     \guidancetrue  ->  \guidancefalse  below. Every box disappears and
%     you get a normal-looking paper.
%  4. \TODO{...} marks things nobody has decided yet. Search for "TODO"
%     before you call any draft finished. \FILL{...} marks a number or
%     name you are expected to substitute.
%=======================================================================

\documentclass[11pt]{article}

%---------------------------------------------------------------- layout
\usepackage[letterpaper,margin=1in]{geometry}
\usepackage{times}
\usepackage[T1]{fontenc}
\usepackage[utf8]{inputenc}
\setlength{\parskip}{0.35em}

%---------------------------------------------------------------- math
\usepackage{amsmath,amssymb,amsfonts}
\usepackage{bm}
\usepackage{siunitx}

%---------------------------------------------------------------- floats
\usepackage{graphicx}
\usepackage{booktabs}
\usepackage{longtable}
\usepackage{array}
\usepackage{multirow}
\usepackage{caption}
\usepackage{subcaption}
\usepackage{pifont}

%---------------------------------------------------------------- code / algos
\usepackage{algorithm}
\usepackage{algpseudocode}
\usepackage{listings}
\lstset{basicstyle=\ttfamily\footnotesize,breaklines=true,
        frame=single,rulecolor=\color{gray!50},columns=fullflexible}

%---------------------------------------------------------------- boxes
\usepackage{xcolor}
\usepackage[most]{tcolorbox}
\usepackage{environ}

%---------------------------------------------------------------- refs
\usepackage[numbers,sort&compress]{natbib}
\bibliographystyle{unsrtnat}
\usepackage[colorlinks=true,linkcolor=blue!60!black,citecolor=blue!60!black,
            urlcolor=blue!60!black]{hyperref}
\usepackage[capitalise,noabbrev]{cleveref}

%=======================================================================
%  GUIDANCE TOGGLE  --  flip to \guidancefalse for a clean draft
%=======================================================================
\newif\ifguidance
\guidancetrue
%\guidancefalse

\NewEnviron{guide}[1][Guidance]{%
  \ifguidance
  \begin{tcolorbox}[breakable,enhanced,colback=blue!4,colframe=blue!55!black,
      boxrule=0.6pt,left=6pt,right=6pt,top=5pt,bottom=5pt,
      title={\footnotesize\bfseries #1},fonttitle=\bfseries,
      coltitle=white,colbacktitle=blue!55!black]
    \footnotesize\BODY
  \end{tcolorbox}
  \fi
}

\NewEnviron{yourturn}[1][Your turn: write this yourself]{%
  \ifguidance
  \begin{tcolorbox}[breakable,enhanced,colback=orange!5,colframe=orange!75!black,
      boxrule=0.6pt,left=6pt,right=6pt,top=5pt,bottom=5pt,
      title={\footnotesize\bfseries #1},fonttitle=\bfseries,
      coltitle=white,colbacktitle=orange!75!black]
    \footnotesize\BODY
  \end{tcolorbox}
  \fi
}

\NewEnviron{figrec}[1][Recommended figure]{%
  \ifguidance
  \begin{tcolorbox}[breakable,enhanced,colback=green!4,colframe=green!45!black,
      boxrule=0.6pt,left=6pt,right=6pt,top=5pt,bottom=5pt,
      title={\footnotesize\bfseries #1},fonttitle=\bfseries,
      coltitle=white,colbacktitle=green!45!black]
    \footnotesize\BODY
  \end{tcolorbox}
  \fi
}

\newcommand{\TODO}[1]{\textcolor{red}{\textbf{[TODO: #1]}}}
\newcommand{\FILL}[1]{\textcolor{blue!65!black}{\underline{\textsf{#1}}}}

%---------------------------------------------------------------- status glyphs
\newcommand{\done}{\textcolor{green!55!black}{\ding{51}}}
\newcommand{\wip}{\textcolor{orange!90!black}{\ding{110}}}
\newcommand{\blocked}{\textcolor{red!80!black}{\ding{55}}}
\newcommand{\notstarted}{\textcolor{gray!70}{\ding{109}}}

%=======================================================================
%  NOTATION MACROS  --  use these everywhere; never hand-type the symbols
%=======================================================================
\newcommand{\pbind}{p_{\mathrm{bind}}}          % binder probability
\newcommand{\yhat}{\hat{y}}                     % Boltz-2 log10 IC50 (uM)
\newcommand{\ICfifty}{\mathrm{IC}_{50}}
\newcommand{\pICfifty}{\mathrm{pIC}_{50}}
\newcommand{\Gmodel}{G}                         % generative model
\newcommand{\Actives}{\mathcal{A}}              % known-actives benchmark set
\newcommand{\Batch}{\mathcal{B}}                % raw generated batch
\newcommand{\Valid}{\mathcal{V}}                % valid, unique molecules
\newcommand{\Passed}{\mathcal{P}}               % molecules passing the gate
\newcommand{\TopK}{\mathcal{K}}                 % top-K selected set
\newcommand{\fpr}{\varphi}                      % fingerprint map
\newcommand{\Tan}{T}                            % Tanimoto coefficient
\newcommand{\simmax}{\mathrm{sim}_{\max}}
\newcommand{\thetahit}{\theta_{\mathrm{hit}}}
\newcommand{\Tbar}{\bar{T}}
\newcommand{\Tstar}{T^{*}}
\newcommand{\Murcko}{\mathrm{Mur}}
\newcommand{\SAscore}{\mathrm{SA}}
\newcommand{\QEDs}{\mathrm{QED}}
\newcommand{\indic}{\mathbb{1}}
\newcommand{\Ngen}{N}
\newcommand{\Kk}{K}

%=======================================================================
\begin{document}

\title{\bfseries An Iterative Closed-Loop Generative Pipeline for De Novo
and Transfer Learning-Guided Small-Molecule Inhibitor Discovery}

\author{%
  \FILL{Author One}$^{1}$, \FILL{Author Two}$^{1}$, \FILL{Author Three}$^{1}$,
  Ibrahim Wichka$^{1}$, \FILL{PI Name}$^{1,*}$ \\[0.4em]
  \small $^{1}$Stevens Institute of Technology, Hoboken, NJ, USA \\
  \small $^{*}$Correspondence: \FILL{pi@stevens.edu}
}
\date{\today}
\maketitle

%-----------------------------------------------------------------------
\begin{abstract}
\noindent
\begin{guide}[Write the abstract LAST]
Target 200--250 words. Use this six-sentence skeleton and expand each into
1--2 sentences:
\begin{enumerate}\itemsep0pt
  \item \textbf{Problem.} Chemical space is enormous; experimental screening
        is slow and expensive.
  \item \textbf{Gap.} Generative models produce molecules, but iterative
        target-specific optimisation and honest validation remain open.
  \item \textbf{What we built.} A closed loop: REINVENT4 generation
        $\rightarrow$ Boltz-2 structure-based scoring $\rightarrow$ gated
        top-$K$ selection $\rightarrow$ transfer-learning fine-tune.
  \item \textbf{How we validated it.} De novo rediscovery: the known-actives
        set is withheld from the generator and used only to measure
        convergence via Tanimoto similarity.
  \item \textbf{Headline result.} Numbers. Actual numbers. $\Tbar_t$ rose
        from \FILL{x} to \FILL{y} over \FILL{n} iterations for
        \FILL{target}; \FILL{n} rediscovery hits at $\thetahit = 0.85$.
  \item \textbf{Why it matters.} Generalisation across two binding-site
        architectures (p38 MAPK14, CDK2) plus an interpretable chemical
        narrative of what the model learns to favour.
\end{enumerate}
Do not put citations in the abstract. Do not use the word ``novel''.
\end{guide}

\TODO{abstract}
\end{abstract}

\noindent\textbf{Keywords:} generative molecular design; de novo drug design;
transfer learning; small-molecule inhibitor discovery; protein--ligand
interaction modeling; structure-based scoring.

%=======================================================================
\section{Introduction}
\label{sec:intro}
%=======================================================================

\begin{guide}[How the Introduction is structured, and why]
An introduction is an argument, not a summary of everything you read. The
argument has a fixed shape and it must end where your paper begins:

\smallskip
\emph{Drug discovery needs better search} $\rightarrow$ \emph{generative models
search well but can't tell good from bad} $\rightarrow$ \emph{structure-based
affinity models can tell good from bad but don't search} $\rightarrow$
\emph{closing the loop between them is the obvious move, and here is what is
still unresolved about doing it} $\rightarrow$ \emph{our contribution.}

\smallskip
Each subsection below is one move in that argument. Write each as
2--4 paragraphs. Delete the \texttt{\textbackslash subsection} headings before submission if
the venue prefers a continuous introduction --- the headings are scaffolding
for you, not necessarily for the reader.

\smallskip
\textbf{Total citation budget:} roughly 35--55 references for the whole
introduction. If a subsection below has fewer than 4 citations, you have not
done the literature review for it yet.
\end{guide}

%-----------------------------------------------------------------------
\subsection{The scale of the search problem}
\label{sec:intro:problem}

\begin{yourturn}[Write \S\ref{sec:intro:problem}: 2 paragraphs]
Establish the stakes with concrete numbers, not adjectives.
\begin{itemize}\itemsep0pt
  \item Size of drug-like chemical space (the $10^{23}$--$10^{60}$ estimates ---
        cite the actual source, do not cite a blog post).
  \item Cost and timeline of bringing one compound to market; attrition rates
        by phase.
  \item What high-throughput screening can and cannot cover, and why "more
        screening" does not close the gap.
\end{itemize}
\textbf{Literature to find:} Polishchuk et al.\ on GDB-17 / enumerable chemical
space; a recent (2020+) R\&D cost-and-attrition analysis; a review of HTS hit
rates. \textbf{Search terms:} \texttt{"size of chemical space drug-like"},
\texttt{"pharmaceutical R\&D productivity attrition"},
\texttt{"high-throughput screening hit rate review"}.
\end{yourturn}

\TODO{write \S\ref{sec:intro:problem}}

%-----------------------------------------------------------------------
\subsection{Generative models for molecular design}
\label{sec:intro:genmodels}

\begin{yourturn}[Write \S\ref{sec:intro:genmodels}: 3--4 paragraphs]
This is your main methods-lineage review. Organise it by \emph{representation
and objective}, not chronologically by year.
\begin{itemize}\itemsep0pt
  \item \textbf{Representations:} SMILES-based RNNs/transformers, graph-based
        generators, 3D/structure-conditioned generators. One paragraph each is
        plenty; you only need enough to justify why a SMILES-based prior is a
        reasonable choice here.
  \item \textbf{Optimisation strategies:} reinforcement learning
        (REINVENT lineage \citep{olivecrona2017, blaschke2020, loeffler2024}),
        transfer learning / fine-tuning on curated corpora
        \citep{segler2018}, genetic algorithms, Bayesian optimisation,
        and active learning.
  \item \textbf{The honest limitation:} generated molecules are only as good as
        the reward. State this plainly --- it is the hinge that sets up
        \S\ref{sec:intro:scoring}.
  \item \textbf{Benchmarking:} GuacaMol, MOSES, and the critique that
        benchmark-optimal molecules are frequently not chemically sensible.
        This motivates our \S\ref{sec:methods:chem} chemical analysis layer.
\end{itemize}
\textbf{Search terms:} \texttt{"de novo molecular generation review 2024"},
\texttt{"SMILES recurrent neural network molecule generation"},
\texttt{"GuacaMol MOSES benchmark generative chemistry"},
\texttt{"reward hacking molecular generative model"}.
\end{yourturn}

\TODO{write \S\ref{sec:intro:genmodels}}

%-----------------------------------------------------------------------
\subsection{Structure-based scoring and binding affinity prediction}
\label{sec:intro:scoring}

\begin{yourturn}[Write \S\ref{sec:intro:scoring}: 3 paragraphs]
Build the case for why Boltz-2 is the scoring function, and be fair about the
alternatives you did not use.
\begin{itemize}\itemsep0pt
  \item \textbf{Classical docking} (AutoDock Vina, Glide, Gold): fast, but
        scoring-function accuracy for affinity ranking is famously weak. Find a
        benchmark paper that quantifies this rather than asserting it.
  \item \textbf{ML scoring functions and FEP:} where each sits on the
        accuracy/cost curve. FEP is accurate and far too slow for
        $5{,}000$ molecules per iteration --- say so with a number.
  \item \textbf{Co-folding foundation models:} AlphaFold2/3
        \citep{jumper2021, abramson2024}, Boltz-1 \citep{wohlwend2024},
        Boltz-2 \citep{boltz2}. Emphasise what makes Boltz-2 usable here
        specifically: it emits a binder probability \emph{and} a regression
        affinity estimate in one call, at a cost that fits an iterative loop.
  \item \textbf{Caveats you must state:} predicted affinity is not measured
        affinity; performance depends on the quality of the input structure;
        the model may be susceptible to systematic exploitation by a generator
        optimising against it.
\end{itemize}
\textbf{Search terms:} \texttt{"docking scoring function affinity ranking
benchmark"}, \texttt{"machine learning scoring function binding affinity
review"}, \texttt{"Boltz-2 binding affinity"}, \texttt{"co-folding model
protein ligand"}.
\end{yourturn}

\TODO{write \S\ref{sec:intro:scoring}}

%-----------------------------------------------------------------------
\subsection{Closed-loop and active-learning pipelines}
\label{sec:intro:closedloop}

\begin{yourturn}[Write \S\ref{sec:intro:closedloop}: 2--3 paragraphs]
This is the subsection that positions \emph{our} contribution, so it needs the
most careful reading. Find prior work that couples a generator to a
structure-based oracle in a loop and describe what each did and where it stopped.
\begin{itemize}\itemsep0pt
  \item Docking-in-the-loop generative pipelines; "generate--dock--retrain"
        systems; DiffDock/AlphaFold-scored generation.
  \item Self-driving / autonomous lab loops, for contrast: they close the loop
        on \emph{experimental} measurements, we close it on predicted ones.
        Be explicit that this is a weaker but far cheaper loop.
  \item \textbf{The gap you are filling.} State it in one sentence you could
        defend in a talk. Candidate framings: (i) most loops report enrichment
        of their own scoring function, which is circular; we validate against a
        withheld set of experimentally measured actives; (ii) most report
        convergence metrics without characterising \emph{what} the chemistry
        did; we track scaffolds, functional groups and synthesisability
        alongside.
\end{itemize}
\textbf{Search terms:} \texttt{"closed loop generative docking active learning
drug design"}, \texttt{"iterative fine-tuning generative model docking score"},
\texttt{"self-driving laboratory molecular design"}.
\end{yourturn}

\TODO{write \S\ref{sec:intro:closedloop}}

%-----------------------------------------------------------------------
\subsection{Target selection: p38 MAPK14 and CDK2}
\label{sec:intro:targets}

\begin{yourturn}[Write \S\ref{sec:intro:targets}: 2 paragraphs, one per target]
Justify both targets on \emph{three} axes each: biology, data availability, and
structural suitability for the scoring model.
\begin{itemize}\itemsep0pt
  \item \textbf{p38 MAPK14:} serine/threonine kinase in inflammatory signalling;
        decades of medicinal chemistry; deep, well-defined, hydrophobic
        ATP site; a large and chemotype-diverse ChEMBL inhibitor set; heavily
        used as a benchmark in prior computational papers, which makes our
        results comparable.
  \item \textbf{CDK2:} cell-cycle kinase; different pocket character and a
        different (smaller) inhibitor landscape; serves as the
        generalisation test.
  \item \textbf{Say what this pair does and does not test.} Two kinases test
        transferability across binding-site architecture but \emph{not} beyond
        the kinase chemotype. Flag this honestly here; return to it in
        \cref{sec:discussion} as a limitation and future direction.
\end{itemize}
\textbf{Literature to find:} 2--3 reviews of p38 inhibitor development
(including why clinical programs failed --- selectivity/toxicity --- which is
good framing); 2--3 on CDK2 inhibitors; the PDB entries you actually used, cited
by their primary papers.
\end{yourturn}

\TODO{write \S\ref{sec:intro:targets}}

%-----------------------------------------------------------------------
\subsection{Rediscovery as validation}
\label{sec:intro:rediscovery}

\begin{yourturn}[Write \S\ref{sec:intro:rediscovery}: 1--2 paragraphs]
Explain and defend the experimental design, because a sceptical reviewer will
push on exactly this.
\begin{itemize}\itemsep0pt
  \item The known-actives set is ground truth and is \emph{never} shown to the
        generator --- not to initialise it, not in the fine-tuning corpus. The
        only place it touches the pipeline is threshold calibration
        (\cref{sec:methods:calib}), and even there it only sets a scalar.
  \item Why this is deliberately conservative: a pipeline that recovers known
        actives without prior structural knowledge of them is evidence that the
        scoring and optimisation components track real structure--activity
        relationships rather than scoring artefacts.
  \item Anticipate the objection: Boltz-2 was trained on public structural and
        affinity data that may include these very compounds. Address it
        directly. Suggested framing: this is a limitation of \emph{any}
        pretrained oracle; it bounds our claim to "the loop navigates toward
        chemistry the oracle recognises as active", and we quantify how far
        that goes rather than assuming it.
\end{itemize}
\end{yourturn}

\TODO{write \S\ref{sec:intro:rediscovery}}

%-----------------------------------------------------------------------
\subsection{Contributions}
\label{sec:intro:contrib}

\begin{guide}[Keep this to 4--5 bullets, each a claim you can point to a figure for]
Every bullet should map onto a results subsection. If a bullet has no
corresponding figure or table, either cut it or go make the figure.
\end{guide}

\noindent In this work we:
\begin{itemize}\itemsep2pt
  \item present a closed-loop pipeline coupling REINVENT4 generation with
        Boltz-2 structure-based scoring through a gated selection function and
        iterative transfer learning; \TODO{confirm phrasing}
  \item calibrate the binder-probability gate $\tau$ empirically against the
        full known-actives benchmark for each target, rather than assuming a
        threshold; \TODO{}
  \item validate the loop as a de novo rediscovery benchmark against
        \FILL{n} withheld experimentally validated inhibitors across two
        kinase targets; \TODO{}
  \item characterise the chemical trajectory of optimisation --- scaffold
        evolution, functional-group enrichment, drug-likeness and synthetic
        accessibility --- across iterations; \TODO{}
  \item release all configurations, SLURM scripts and analysis notebooks
        \TODO{repository URL and license}.
\end{itemize}

\begin{figrec}[Figure 1 --- pipeline schematic (make this early, it anchors everything)]
A single-column block diagram of the closed loop:
REINVENT4 $\Gmodel_t$ $\rightarrow$ Boltz-2 scoring $\rightarrow$ gating
function $\rightarrow$ top-$\Kk$ $\rightarrow$ transfer learning $\rightarrow$
$\Gmodel_{t+1}$, with the known-actives set drawn \emph{off} the main path
(dashed) feeding only calibration and the convergence check.
\textbf{Source:} adapt the flowchart already in the research plan.
\textbf{Tool:} TikZ, or draw.io exported to PDF. Do not use a raster screenshot.
\end{figrec}

% \begin{figure}[t]
%   \centering
%   \includegraphics[width=\linewidth]{figures/fig1_pipeline_schematic.pdf}
%   \caption{Overview of the closed-loop pipeline. \TODO{caption}}
%   \label{fig:pipeline}
% \end{figure}

%=======================================================================
\section{Methods}
\label{sec:methods}
%=======================================================================

\begin{guide}[How to use the Methods section]
Everything in this section that is \emph{shared infrastructure} --- the loop
structure, the notation, the scoring maths, the metric definitions --- is
already written for you. Do not rewrite it; you would only introduce
inconsistencies between sections written by different people.

\smallskip
What you \emph{do} write is inside the orange boxes: the decisions you made,
the numbers you chose, the things that went wrong and what you did about them.
Write those in past tense, first person plural ("we set", "we observed"),
and be specific. "We used a GPU" is not a method; "each Boltz-2 call ran on one
NVIDIA A100 with \texttt{--diffusion\_samples 1} and
\texttt{--sampling\_steps 200}" is.

\smallskip
\textbf{Use the notation macros.} \texttt{\textbackslash pbind}, \texttt{\textbackslash Tbar}, \texttt{\textbackslash TopK},
etc.\ are defined in the preamble. If you hand-type $p_{bind}$ somewhere it
will render differently from everywhere else and a reader will notice.
\end{guide}

%-----------------------------------------------------------------------
\subsection{Overview of the closed loop}
\label{sec:methods:overview}

At iteration $t$, a generative model $\Gmodel_t$ samples a batch of
$\Ngen$ molecules. Each valid molecule is scored against the prepared target
structure by Boltz-2, which returns a binder probability and a regression-based
potency estimate. A gated scoring function discards molecules below a
binder-probability threshold $\tau$ and ranks the remainder by predicted
potency; the top $\Kk$ molecules form a training corpus used to fine-tune
$\Gmodel_t$ by transfer learning, producing $\Gmodel_{t+1}$. Progress is
measured against a held-out set $\Actives$ of experimentally validated
inhibitors that is never exposed to the generator.
\Cref{alg:loop} states the procedure.

\begin{algorithm}[t]
\caption{Closed-loop generative optimisation for one target}
\label{alg:loop}
\begin{algorithmic}[1]
\Require pretrained prior $\Gmodel_1$; target structure $S$; gate $\tau$;
         batch size $\Ngen$; corpus size $\Kk$; hit threshold $\thetahit$;
         held-out actives $\Actives$
\For{$t = 1, 2, \dots, T_{\max}$}
  \State $\Batch_t \gets \{m_1,\dots,m_\Ngen\} \sim \Gmodel_t$
         \Comment{sample SMILES}
  \State $\Valid_t \gets \mathrm{dedup}(\mathrm{sanitise}(\Batch_t))$
         \Comment{canonicalise, valence check, drop duplicates}
  \ForAll{$m \in \Valid_t$}
    \State $(\pbind(m),\, \yhat(m)) \gets \textsc{Boltz-2}(S, m)$
  \EndFor
  \State $\Passed_t \gets \{\, m \in \Valid_t : \pbind(m) \ge \tau \,\}$
  \State $\TopK_t \gets \operatorname*{arg\,top-}\Kk_{\,m \in \Passed_t} S(m)$
         \Comment{\cref{eq:score}}
  \State compute $\Tbar_t,\ \Tstar_t$ against $\Actives$
         \Comment{\cref{eq:tbar,eq:tstar}}
  \If{$\Tstar_t \ge \thetahit$}
    \State \textbf{record} rediscovery hit at iteration $t$
  \EndIf
  \State $\Gmodel_{t+1} \gets \textsc{TransferLearn}(\Gmodel_t, \TopK_t)$
\EndFor
\end{algorithmic}
\end{algorithm}

%-----------------------------------------------------------------------
\subsection{Datasets and target preparation}
\label{sec:methods:data}

\subsubsection{Known-actives benchmark sets}

For each target we assembled a set $\Actives$ of experimentally validated
small-molecule inhibitors with measured $\ICfifty$ values, retrieved from
ChEMBL \citep{zdrazil2024}. Records were filtered to $\ICfifty$ endpoints,
deduplicated, and canonicalised as SMILES. Measured potencies are expressed
throughout as
\begin{equation}
  \pICfifty = -\log_{10}\!\left(\ICfifty\,[\mathrm{M}]\right).
  \label{eq:pic50}
\end{equation}
Set sizes after curation are given in \cref{tab:datasets}. This set is used
exclusively (i) to calibrate the gate $\tau$ (\cref{sec:methods:calib}) and
(ii) as the convergence reference (\cref{sec:methods:convergence}). It is never
used to initialise, condition, or fine-tune the generative model.

\begin{table}[t]
\centering
\caption{Known-actives benchmark sets and prepared target structures.}
\label{tab:datasets}
\small
\begin{tabular}{lccccc}
\toprule
Target & UniProt & PDB ID & Resolution (\si{\angstrom}) & $|\Actives|$ & Median $\pICfifty$ \\
\midrule
p38 MAPK14 & \FILL{Q16539} & \FILL{PDB} & \FILL{x.xx} & 3{,}860 & \FILL{x.xx} \\
CDK2       & \FILL{P24941} & \FILL{PDB} & \FILL{x.xx} & 1{,}402 & \FILL{x.xx} \\
\bottomrule
\end{tabular}
\end{table}

\begin{yourturn}[Fill \cref{tab:datasets} and write 1 paragraph on curation]
The two $|\Actives|$ counts above are the manifest row counts actually docked
(\texttt{threshold/<TARGET>/manifest.csv}). Verify them against your own
\texttt{prepare\_inputs.py} output before submission and update if they differ.
Then write a short paragraph covering:
\begin{itemize}\itemsep0pt
  \item the exact ChEMBL query / export settings and the download date;
  \item what you filtered out and why (assay type, censored values such as
        \texttt{>} or \texttt{<}, salts/mixtures, missing SMILES,
        unparseable SMILES);
  \item how many records went in and how many survived --- give both numbers;
  \item whether duplicate measurements of the same compound were averaged or
        the highest-confidence record kept.
\end{itemize}
\end{yourturn}

\subsubsection{Protein structure preparation}

\begin{yourturn}[Write this: 1 paragraph per target]
State for each target: the PDB entry chosen and why (resolution,
co-crystallised ligand, conformational state --- DFG-in/DFG-out matters for
p38); crystallographic waters and heteroatoms removed; hydrogens added;
protonation state assignment and at what pH; how the binding site was defined
for the Boltz-2 input; and the exact protein sequence used in the YAML
(state whether it is the full construct or a truncation, and give residue
ranges). If you used a predicted structure anywhere, say so and note the
caveat that scoring quality is contingent on structural accuracy
\citep{abramson2024}.
\end{yourturn}

\TODO{structure preparation paragraphs}

\subsubsection{Molecular representation}

All generated and reference molecules are canonicalised as SMILES using
RDKit \citep{rdkit}. For similarity computations, each molecule $m$ is mapped
to a Morgan (ECFP4-equivalent) fingerprint \citep{rogers2010}
\begin{equation}
  \fpr(m) \in \{0,1\}^{2048},
  \qquad \text{radius } r = 2, \ \text{fold size } 2048,
  \label{eq:fp}
\end{equation}
and the Tanimoto coefficient between two molecules is
\begin{equation}
  \Tan\!\left(\fpr(m), \fpr(m')\right)
  = \frac{\left|\fpr(m) \wedge \fpr(m')\right|}
         {\left|\fpr(m) \vee \fpr(m')\right|} \in [0,1].
  \label{eq:tanimoto}
\end{equation}

%-----------------------------------------------------------------------
\subsection{Molecular generation}
\label{sec:methods:gen}

Generation uses REINVENT4 \citep{loeffler2024} in sampling mode
(\texttt{run\_type = "sampling"}), which draws from the model without
modifying its weights. At $t=1$ the model $\Gmodel_1$ is the unmodified
pretrained prior \texttt{reinvent\_pubchem.prior}, trained on PubChem and
carrying no knowledge of either target; for $t \ge 2$, $\Gmodel_t$ is the
checkpoint produced by the transfer-learning update of iteration $t-1$
(\cref{sec:methods:tl}). Iteration 1 is therefore the only round sampled from a
fully unbiased distribution.

Each iteration samples $\Ngen = 5{,}000$ SMILES per target with
\texttt{unique\_molecules = true} (post-hoc deduplication) and
\texttt{randomize\_smiles = true} (atom-ordering augmentation during sampling).
The batch size is a fixed design choice applied uniformly to both targets and
is unrelated to the size of either target's benchmark set. Sampled strings are
sanitised and canonicalised with RDKit; molecules failing valence or
sanitisation checks are discarded, so $|\Valid_t| \le \Ngen$.

\begin{yourturn}[Write 1 short paragraph + fill \cref{tab:genconfig}]
Report, per target and per iteration: how many of the $5{,}000$ sampled strings
survived canonicalisation and deduplication, and whether that fraction changed
as iterations progressed. A \emph{falling} valid fraction across iterations is
worth commenting on --- it can indicate the fine-tuned model drifting toward
degenerate strings. Also state the software version
(\texttt{reinvent --version}) and the random seed if you set one.
\end{yourturn}

\begin{table}[h]
\centering
\caption{REINVENT4 sampling configuration (identical across targets and iterations).}
\label{tab:genconfig}
\small
\begin{tabular}{lll}
\toprule
Parameter & Value & Note \\
\midrule
\texttt{run\_type}          & \texttt{sampling} & weights unchanged \\
\texttt{model\_file}        & prior ($t=1$) / fine-tuned ckpt ($t\ge2$) & \cref{sec:methods:tl} \\
\texttt{num\_smiles}        & 5{,}000 & fixed by design, both targets \\
\texttt{unique\_molecules}  & \texttt{true} & post-hoc deduplication \\
\texttt{randomize\_smiles}  & \texttt{true} & atom-order augmentation \\
\texttt{device}             & \texttt{cuda:0} & one NVIDIA A100 \\
REINVENT4 version           & \FILL{x.y.z} & \\
\bottomrule
\end{tabular}
\end{table}

%-----------------------------------------------------------------------
\subsection{Structure-based scoring with Boltz-2}
\label{sec:methods:boltz}

Every molecule is evaluated by a single Boltz-2 \citep{boltz2} call taking the
target protein sequence and one ligand SMILES, with an explicit
\texttt{properties: affinity: binder} declaration in the input YAML so that the
affinity head is invoked. Each call returns:
\begin{itemize}\itemsep0pt
  \item a predicted complex structure and confidence metrics
        (pTM, ipTM, pLDDT);
  \item $\pbind(m) \in [0,1]$, the predicted probability that $m$ is a binder
        at the site (\texttt{affinity\_probability\_binary});
  \item $\yhat(m) = \log_{10}\widehat{\ICfifty}(m)$ with
        $\widehat{\ICfifty}$ in \textmu M
        (\texttt{affinity\_pred\_value}); lower is more potent.
\end{itemize}
For comparison with experimental values on a common scale we convert
\begin{equation}
  \widehat{\pICfifty}(m) \;=\; 6 - \yhat(m),
  \label{eq:convert}
\end{equation}
which follows from \cref{eq:pic50} and the \textmu M unit convention.

Identical inference settings were used for calibration and for scoring
generated molecules:
\texttt{--model boltz2},
\texttt{--diffusion\_samples 1},
\texttt{--sampling\_steps 200},
\texttt{--diffusion\_samples\_affinity 5},
\texttt{--affinity\_mw\_correction}.

\begin{yourturn}[Write 1 paragraph]
Cover: Boltz-2 version/commit; wall-clock cost per molecule and total GPU-hours
consumed per target; how failures were handled (molecules that produced no
affinity JSON and were dropped, and how many); whether the calibration array
job and the generated-molecule docking job used the same code path (they should
--- if they differed at all, say how). If you hit the \SI{24}{\hour} walltime on
the generated-molecule docking run, describe what you did and whether any
molecules were re-docked.
\end{yourturn}

%-----------------------------------------------------------------------
\subsection{Threshold calibration}
\label{sec:methods:calib}

\begin{guide}[Why this subsection exists]
$\tau$ is the one number in the whole pipeline chosen by looking at the actives.
Reviewers will ask whether that leaks information. Your defence is that $\tau$
is a scalar derived from a \emph{score distribution}, not from any structural
feature of any active, and that it is fixed before iterative optimisation
begins. Make sure the text says this explicitly.
\end{guide}

Before iterative optimisation, we characterised the Boltz-2 scoring landscape
for each target. We evaluated (i) the complete known-actives set $\Actives$ as a
positive control and (ii) a background sample of
$n_{\mathrm{bg}} = \FILL{500\text{--}1{,}000}$ molecules drawn from the
unmodified prior with no reward signal.

Two properties were required before committing a target to iterative runs.
First, \emph{separation}: the actives must score above the background on both
outputs. We quantify this with the rank-biserial form of the Mann--Whitney $U$
statistic, equivalent to the probability that a randomly chosen active
out-scores a randomly chosen background molecule,
\begin{equation}
  \mathrm{AUC}
  = \Pr\!\left[\pbind(a) > \pbind(b)\right],
  \qquad a \sim \Actives,\; b \sim \mathrm{background},
  \label{eq:auc}
\end{equation}
with $\mathrm{AUC} = 0.5$ indicating no discrimination. Second,
\emph{potency correlation}: predicted versus experimental potency across
$\Actives$, reported as Pearson $r$ and Spearman $\rho$ between
$\widehat{\pICfifty}$ (\cref{eq:convert}) and measured $\pICfifty$.

The gate $\tau$ is then chosen to sit meaningfully above the background median
while remaining attainable by a non-trivial fraction of unbiased samples:
\begin{equation}
  \tau \;=\;
  \operatorname*{arg\,max}_{\tau' \in \{0.3,\,0.4,\,0.5\}}
  \Big\{\,\tau' \;:\;
    \Pr_{b}\!\left[\pbind(b) \ge \tau'\right] \ge \epsilon
  \,\Big\},
  \label{eq:tau}
\end{equation}
where $\epsilon = \FILL{0.0x}$ is the minimum background pass rate needed to
keep the transfer-learning corpus populated in early iterations. A gate that
excludes every unbiased molecule suppresses the learning signal entirely; a
gate that admits the majority pollutes the corpus with non-binders.

\begin{yourturn}[Fill \cref{tab:calibration} and write 1--2 paragraphs]
Report per target: the $\pbind$ percentiles (10/25/50/75/90) for both actives
and background; the fraction of each population clearing $0.3$, $0.4$, $0.5$;
the AUC from \cref{eq:auc}; Pearson $r$ and Spearman $\rho$; and the $\tau$ you
finally chose \emph{with the sentence of reasoning that led to it}. If the two
targets ended up with different $\tau$, say so and justify it --- that is a
legitimate finding, not an inconsistency to hide. Also fix $\epsilon$ in
\cref{eq:tau} to the value you actually used, or replace the equation with a
plain sentence if you selected $\tau$ by inspection.
\end{yourturn}

\begin{table}[h]
\centering
\caption{Boltz-2 calibration summary. Percentiles are of $\pbind$.}
\label{tab:calibration}
\small
\begin{tabular}{llcccccc}
\toprule
Target & Population & $P_{10}$ & $P_{25}$ & $P_{50}$ & $P_{75}$ & $P_{90}$ & $\Pr[\pbind \ge \tau]$ \\
\midrule
\multirow{2}{*}{p38 MAPK14}
  & known actives & \FILL{} & \FILL{} & \FILL{} & \FILL{} & \FILL{} & \FILL{} \\
  & prior background & \FILL{} & \FILL{} & \FILL{} & \FILL{} & \FILL{} & \FILL{} \\
\midrule
\multirow{2}{*}{CDK2}
  & known actives & \FILL{} & \FILL{} & \FILL{} & \FILL{} & \FILL{} & \FILL{} \\
  & prior background & \FILL{} & \FILL{} & \FILL{} & \FILL{} & \FILL{} & \FILL{} \\
\bottomrule
\end{tabular}
\end{table}

%-----------------------------------------------------------------------
\subsection{Gated scoring function and candidate selection}
\label{sec:methods:score}

The optimisation objective is a gated score. For molecule $m$,
\begin{equation}
  S(m) =
  \begin{cases}
    -\,\yhat(m), & \pbind(m) \ge \tau, \\[2pt]
    0,           & \text{otherwise},
  \end{cases}
  \label{eq:score}
\end{equation}
where the negation makes lower predicted $\ICfifty$ (higher potency)
correspond to a higher score, consistent with maximisation. Molecules below the
gate are excluded from all downstream steps. Among the survivors
$\Passed_t = \{m \in \Valid_t : \pbind(m) \ge \tau\}$, candidates are ranked by
$S(m)$ and the top $\Kk$ retained:
\begin{equation}
  \TopK_t
  = \operatorname*{arg\,top-}{\Kk}_{\,m \in \Passed_t} S(m),
  \qquad \Kk \in [200, 500].
  \label{eq:topk}
\end{equation}
When $|\Passed_t| < \Kk$ the entire passing set is used.

\begin{guide}[Alternative scoring formulations --- describe these only if you ran them]
The plan lists two alternatives. If you ran either as an ablation, report it in
\cref{sec:results:ablation}; if not, mention them in one sentence in
\cref{sec:discussion} as untested design choices rather than claiming them.
\begin{itemize}\itemsep0pt
  \item \textbf{Option A, composite:}
        $S_A(m) = w_1\,\pbind(m) + w_2\,(-\yhat(m))$ --- a soft gradient over
        both outputs, at the risk of admitting high-$\pbind$ non-binders.
  \item \textbf{Option C, hard filter then potency only:} $\pbind$ used purely
        as a binary filter, ranking by $\yhat$ alone --- a cleaner potency
        signal that discards the probabilistic gradient which may be
        informative when few molecules pass.
\end{itemize}
\end{guide}

\begin{yourturn}[State your actual $\tau$ and $\Kk$, and why]
One short paragraph. Give the final values used for each target. If $\Kk$ was
adapted because too few molecules passed the gate in early iterations, describe
the rule you used --- an unstated adaptive $\Kk$ is a reproducibility hole.
\end{yourturn}

%-----------------------------------------------------------------------
\subsection{Transfer-learning update}
\label{sec:methods:tl}

The corpus $\TopK_t$ is used to fine-tune $\Gmodel_t$ in REINVENT4 transfer
learning mode (\texttt{run\_type = "transfer\_learning"}), shifting the model's
output distribution toward the chemical features associated with high $S(m)$
and producing the checkpoint $\Gmodel_{t+1}$. Formally, fine-tuning minimises
the negative log-likelihood of the selected corpus under the model,
\begin{equation}
  \mathcal{L}(\theta)
  = -\frac{1}{|\TopK_t|}
    \sum_{m \in \TopK_t} \sum_{i=1}^{|m|}
      \log P_\theta\!\left(s_i \mid s_{<i}\right),
  \label{eq:tlloss}
\end{equation}
where $s_1,\dots,s_{|m|}$ are the SMILES tokens of $m$ and $\theta$ is
initialised from $\Gmodel_t$ rather than from scratch. Each iteration's prior is
therefore the previous iteration's fine-tuned checkpoint, which is what makes
the loop compounding.

\begin{yourturn}[Fill \cref{tab:tlconfig} and write 1 paragraph]
Hyperparameters must be reported --- this is the step that determines whether
the loop converges or collapses. Also address: were these held fixed across all
iterations (they should be, unless an ablation says otherwise)? Did you monitor
validation loss or early-stop? Did you observe mode collapse --- a sharp drop in
unique molecules or unique scaffolds after a fine-tune --- and if so at which
iteration?
\end{yourturn}

\begin{table}[h]
\centering
\caption{Transfer-learning hyperparameters, held fixed across iterations.}
\label{tab:tlconfig}
\small
\begin{tabular}{ll}
\toprule
Parameter & Value \\
\midrule
Corpus size $\Kk$              & \FILL{} \\
Epochs                         & \FILL{} \\
Learning rate                  & \FILL{} \\
Batch size                     & \FILL{} \\
Optimiser                      & \FILL{} \\
Validation split               & \FILL{none / x\%} \\
Early stopping                 & \FILL{} \\
\bottomrule
\end{tabular}
\end{table}

%-----------------------------------------------------------------------
\subsection{Convergence metrics}
\label{sec:methods:convergence}

Convergence is measured by structural similarity between generated molecules
and the held-out actives $\Actives$. For each generated molecule $m$ we define
its nearest-neighbour similarity to the benchmark set,
\begin{equation}
  \simmax(m, \Actives)
  = \max_{a \in \Actives}\,
    \Tan\!\left(\fpr(m), \fpr(a)\right).
  \label{eq:simmax}
\end{equation}
Two summary statistics are recorded per iteration. The \emph{mean
nearest-neighbour similarity} over the selected set,
\begin{equation}
  \Tbar_t
  = \frac{1}{|\TopK_t|} \sum_{m \in \TopK_t} \simmax(m, \Actives),
  \label{eq:tbar}
\end{equation}
is the primary convergence signal: a monotonically increasing $\Tbar_t$
indicates the loop is navigating toward known-active chemical space. The
\emph{maximum} similarity,
\begin{equation}
  \Tstar_t = \max_{m \in \TopK_t}\, \simmax(m, \Actives),
  \label{eq:tstar}
\end{equation}
is used for hit detection: molecule $m$ is declared a \emph{rediscovery hit}
when $\simmax(m,\Actives) \ge \thetahit$, with $\thetahit = 0.85$ for
near-identical scaffolds and $\thetahit = 1.0$ for an exact canonical SMILES
match. The first iteration at which $\Tstar_t \ge \thetahit$ is recorded as the
convergence iteration. As a null reference we also report $\Tbar_0$, computed
identically on an unbiased sample from the prior, so that any rise in $\Tbar_t$
is read against the level expected by chance.

Because the convergence iteration is not known in advance, no fixed stopping
criterion is imposed a priori; $\Tbar_t$ and $\Tstar_t$ are monitored per
iteration and the run terminates on hit detection or on reaching
$T_{\max} = \FILL{n}$ iterations.

\begin{guide}[Two things to compute that the plan does not require but reviewers will want]
\begin{enumerate}\itemsep0pt
  \item \textbf{A random-baseline curve.} Run the metric on $\Kk$ molecules
        sampled from the untouched prior at every iteration. If $\Tbar_t$ rises
        and the baseline is flat, your effect is real. Without it, the rise is
        uninterpretable.
  \item \textbf{Error bars.} $\Tbar_t$ is a mean over $\Kk$ molecules --- report
        the standard error, or better, bootstrap a 95\% CI. A convergence curve
        with no uncertainty band is the most common criticism of this kind of
        plot.
\end{enumerate}
\end{guide}

%-----------------------------------------------------------------------
\subsection{Iterative chemical analysis}
\label{sec:methods:chem}

In parallel with convergence tracking, the selected set $\TopK_t$ is
characterised chemically at each iteration.

\paragraph{Scaffold tracking.}
Each molecule is reduced to its Bemis--Murcko scaffold \citep{bemis1996},
$\Murcko(m)$, using the RDKit implementation. Scaffold diversity at iteration
$t$ is the normalised count of distinct scaffolds,
\begin{equation}
  D_t = \frac{\left|\{\, \Murcko(m) : m \in \TopK_t \,\}\right|}{|\TopK_t|}
  \in (0, 1],
  \label{eq:scaffolddiv}
\end{equation}
with $D_t \to 0$ indicating collapse onto a dominant chemotype. We additionally
report the scaffold overlap with the benchmark set,
$\left|\{\Murcko(m) : m \in \TopK_t\} \cap \{\Murcko(a) : a \in \Actives\}\right|$,
and the scaffold transition between consecutive iterations, comparing the most
frequent scaffolds of $\TopK_t$ and $\TopK_{t+1}$ to identify which ring systems
the model retains, discards, or introduces.

\paragraph{Functional-group trajectories.}
For a set $\mathcal{G}$ of SMARTS patterns covering hydrogen-bond donors and
acceptors (hydroxyl, amine, carbonyl, amide), halogens (F, Cl, Br), aromatic
systems, and basic nitrogen heterocycles (piperidine, piperazine, morpholine),
the frequency of group $g$ at iteration $t$ is
\begin{equation}
  f_{g,t}
  = \frac{1}{|\TopK_t|}
    \sum_{m \in \TopK_t} \indic\!\left[m \text{ matches } g\right].
  \label{eq:fg}
\end{equation}
Monotone trends in $f_{g,t}$ are interpreted against the known character of the
target binding site. The amine pattern excludes amide nitrogens so that amides
are not double-counted.

\paragraph{R-group decomposition.}
Where a dominant core emerges, molecules matching it are decomposed into core
plus substituents at each attachment point, and substituents at position $R_j$
are ranked by mean predicted score,
\begin{equation}
  \bar{S}_{R_j = v}
  = \frac{1}{|\{m : R_j(m) = v\}|}
    \sum_{m\,:\,R_j(m) = v} S(m),
  \label{eq:sar}
\end{equation}
yielding a computational structure--activity table.

\paragraph{Drug-likeness and synthesisability.}
For every molecule in $\TopK_t$ we compute molecular weight, calculated
$\log P$, hydrogen-bond donor and acceptor counts and topological polar surface
area; the fraction violating each Lipinski threshold \citep{lipinski1997}
($\mathrm{MW} > 500$, $c\log P > 5$, $\mathrm{HBD} > 5$, $\mathrm{HBA} > 10$);
QED \citep{bickerton2012}; and the Ertl--Schuffenhauer synthetic accessibility
score \citep{ertl2009}, on a scale from 1 (easy) to 10 (difficult). Together
these define a chemical feasibility envelope; systematic drift outside it across
iterations is flagged as possible over-optimisation against the scoring
function at the expense of chemical viability.

\paragraph{Red-flag screening.}
Candidates are screened against the RDKit PAINS filter catalogue
\citep{baell2010}. A PAINS match is recorded as an annotation for expert review,
not as an automatic disqualification.

\paragraph{Descriptor-space projection.}
A PCA is fitted on the standardised descriptor matrix of $\Actives$ alone
(MW, $c\log P$, HBD, HBA, TPSA, QED) and each iteration's $\TopK_t$ is projected
into that fixed space. Fitting on the actives only keeps the axes constant
across iterations so that successive plots are directly comparable.

\begin{yourturn}[Write 1 paragraph on expert structural annotation]
Describe your manual review protocol from the plan's Section 3.6: how many
molecules per iteration (5--10), how they were selected (diversity-prioritised,
not just top-scoring), who reviewed them, and what was recorded --- structural
novelty relative to $\Actives$ and the literature, qualitative synthetic
feasibility, plausibility of the predicted binding mode, and red-flag
substructures. Say plainly that this step is qualitative.
\end{yourturn}

%-----------------------------------------------------------------------
\subsection{Computational environment and reproducibility}
\label{sec:methods:compute}

All generation, scoring, and fine-tuning jobs ran on the Anvil cluster
(Purdue University / NSF ACCESS) under the SLURM scheduler, using
NVIDIA A100 GPUs. REINVENT4 and Boltz-2 were installed in separate conda
environments (Python \FILL{3.11} and \FILL{3.10} respectively) because of
incompatible PyTorch and CUDA dependencies. Chemical analysis was performed in
Python with RDKit \citep{rdkit}, pandas, scikit-learn, and matplotlib.

\begin{yourturn}[Fill in and complete]
Report software versions for everything in \cref{tab:software}, total GPU-hours
per target, and the repository URL and commit hash corresponding to the results
in this paper. A version table is the cheapest possible reproducibility win ---
do not skip it.
\end{yourturn}

\begin{table}[h]
\centering
\caption{Software versions.}
\label{tab:software}
\small
\begin{tabular}{ll@{\hspace{2em}}ll}
\toprule
Component & Version & Component & Version \\
\midrule
REINVENT4 & \FILL{} & RDKit      & \FILL{} \\
Boltz-2   & \FILL{} & scikit-learn & \FILL{} \\
PyTorch   & \FILL{} & pandas     & \FILL{} \\
CUDA      & \FILL{12.6} & Python & \FILL{} \\
\bottomrule
\end{tabular}
\end{table}

\paragraph{Data and code availability.}
\TODO{repository URL, commit hash, DOI/Zenodo archive, license}

%=======================================================================
\section{Results}
\label{sec:results}
%=======================================================================

\begin{guide}[How to use the Results section --- read this before writing anything]
This section is deliberately built as a \textbf{living tracker first, paper
second}. \Cref{sec:results:tracker} is a status board you update as work lands;
\cref{sec:results:p38,sec:results:cdk2} contain one short log block per
iteration that you fill in the same week the iteration finishes, while you still
remember what happened.

\smallskip
When it is time to write the actual prose Results, you will not be starting from
a blank page --- you will be summarising your own logs. That is the entire point
of the structure. Keep the logs even after the prose exists; move them to an
appendix or delete them at the very end.

\smallskip
\textbf{Rules for the prose you eventually write:}
\begin{itemize}\itemsep0pt
  \item Results state what happened. Interpretation goes in
        \cref{sec:discussion}. "$\Tbar_t$ rose from $0.21$ to $0.34$" belongs
        here; "suggesting the model learned the hinge-binding motif" does not.
  \item Every number in the text must also appear in a table or figure, and
        vice versa.
  \item Report the failures. An iteration where $\Tbar_t$ dropped, a target
        where calibration showed no separation, a fine-tune that collapsed
        diversity --- these make the paper more credible, not less.
\end{itemize}
\end{guide}

%-----------------------------------------------------------------------
\subsection{Pipeline progress tracker}
\label{sec:results:tracker}

\begin{guide}[Legend and instructions]
Update \cref{tab:status-p38,tab:status-cdk2} as each stage completes.
Glyphs: \done\ complete and pushed to \texttt{main}\quad
\wip\ in progress / job running\quad
\blocked\ blocked or failed\quad
\notstarted\ not started.
Add a new row for each iteration as it begins. The "Owner" column is who to ask
when something in that row is unclear six weeks from now.
\end{guide}

\begin{table}[h]
\centering
\caption{Execution status --- p38 MAPK14. Owner initials in parentheses.}
\label{tab:status-p38}
\small
\begin{tabular}{clccccl}
\toprule
Iter & Stage set & Gen & Dock & Gate + TL & Chem.\ analysis & Last update \\
\midrule
0 & calibration ($\tau$) & --- & \wip\ (\FILL{MB}) & --- & \notstarted & \FILL{date} \\
1 & prior sampling       & \done\ (\FILL{}) & \wip\ (\FILL{}) & \notstarted & \notstarted & \FILL{date} \\
2 & \FILL{}              & \notstarted & \notstarted & \notstarted & \notstarted & \FILL{date} \\
3 & \FILL{}              & \notstarted & \notstarted & \notstarted & \notstarted & \FILL{date} \\
\bottomrule
\end{tabular}
\end{table}

\begin{table}[h]
\centering
\caption{Execution status --- CDK2. Owner initials in parentheses.}
\label{tab:status-cdk2}
\small
\begin{tabular}{clccccl}
\toprule
Iter & Stage set & Gen & Dock & Gate + TL & Chem.\ analysis & Last update \\
\midrule
0 & calibration ($\tau$) & --- & \wip\ (\FILL{AK}) & --- & \notstarted & \FILL{date} \\
1 & prior sampling       & \done\ (\FILL{}) & \wip\ (\FILL{}) & \notstarted & \notstarted & \FILL{date} \\
2 & \FILL{}              & \notstarted & \notstarted & \notstarted & \notstarted & \FILL{date} \\
3 & \FILL{}              & \notstarted & \notstarted & \notstarted & \notstarted & \FILL{date} \\
\bottomrule
\end{tabular}
\end{table}

\noindent
\Cref{tab:master} is the quantitative counterpart: one row per completed
iteration per target, and the table that most of the Results prose will be
written from.

\begin{table}[t]
\centering
\caption{Master results tracker. One row per completed iteration.
$|\Valid_t|$: valid unique molecules of $\Ngen = 5{,}000$ sampled.
$|\Passed_t|$: molecules with $\pbind \ge \tau$.
$\bar{S}$: mean gated score over $\TopK_t$.
$\Tbar_t$, $\Tstar_t$: \cref{eq:tbar,eq:tstar}.
$D_t$: scaffold diversity, \cref{eq:scaffolddiv}.
Hits: molecules with $\simmax \ge \thetahit = 0.85$.}
\label{tab:master}
\small
\setlength{\tabcolsep}{4pt}
\begin{tabular}{llrrrrrrrrrr}
\toprule
Target & $t$ & $|\Valid_t|$ & $|\Passed_t|$ & $\Kk$ & $\bar{S}$ &
$\Tbar_t$ & $\Tstar_t$ & Hits & $D_t$ & $\overline{\QEDs}$ & $\overline{\SAscore}$ \\
\midrule
\multirow{4}{*}{p38 MAPK14}
 & 0\textsuperscript{a} & --- & --- & --- & --- & \FILL{} & \FILL{} & \FILL{} & \FILL{} & \FILL{} & \FILL{} \\
 & 1 & \FILL{} & \FILL{} & \FILL{} & \FILL{} & \FILL{} & \FILL{} & \FILL{} & \FILL{} & \FILL{} & \FILL{} \\
 & 2 & \FILL{} & \FILL{} & \FILL{} & \FILL{} & \FILL{} & \FILL{} & \FILL{} & \FILL{} & \FILL{} & \FILL{} \\
 & 3 & \FILL{} & \FILL{} & \FILL{} & \FILL{} & \FILL{} & \FILL{} & \FILL{} & \FILL{} & \FILL{} & \FILL{} \\
\midrule
\multirow{4}{*}{CDK2}
 & 0\textsuperscript{a} & --- & --- & --- & --- & \FILL{} & \FILL{} & \FILL{} & \FILL{} & \FILL{} & \FILL{} \\
 & 1 & \FILL{} & \FILL{} & \FILL{} & \FILL{} & \FILL{} & \FILL{} & \FILL{} & \FILL{} & \FILL{} & \FILL{} \\
 & 2 & \FILL{} & \FILL{} & \FILL{} & \FILL{} & \FILL{} & \FILL{} & \FILL{} & \FILL{} & \FILL{} & \FILL{} \\
 & 3 & \FILL{} & \FILL{} & \FILL{} & \FILL{} & \FILL{} & \FILL{} & \FILL{} & \FILL{} & \FILL{} & \FILL{} \\
\midrule
\multicolumn{12}{l}{\footnotesize
  \textsuperscript{a}Row $t=0$ is the unbiased-prior null reference: the same
  metrics computed on $\Kk$ molecules sampled from} \\
\multicolumn{12}{l}{\footnotesize
  \phantom{a}the untouched prior. Every claim of improvement is relative to this row.} \\
\bottomrule
\end{tabular}
\end{table}

%-----------------------------------------------------------------------
\subsection{Boltz-2 calibration}
\label{sec:results:calib}

\begin{yourturn}[Write 2--3 paragraphs. This is the first real result in the paper.]
Report, per target: whether known actives separated from the prior background
on $\pbind$, by how much (\cref{eq:auc}), and how predicted potency correlated
with measured potency ($r$, $\rho$, $n$). Then state the chosen $\tau$ and the
fraction of background molecules it admits.

\smallskip
\textbf{This subsection carries a go/no-go decision, so be direct about it.}
If a target showed weak or absent separation, say so here and explain what you
concluded --- structural preparation issue, or a genuine limitation of Boltz-2
for that target. Either is a publishable finding. Burying it is not an option,
because a reader can tell from the correlation plot.
\end{yourturn}

\TODO{calibration results}

\begin{figrec}[Figure 2 --- calibration (2 panels)]
\textbf{(a)} $\pbind$ distributions, known actives vs.\ prior background,
overlaid density histograms per target, with vertical lines at
$\tau \in \{0.3, 0.4, 0.5\}$.
\textbf{(b)} Predicted vs.\ experimental $\pICfifty$ scatter, with $y=x$
reference line, Pearson $r$ annotated, and a marginal histogram.
\textbf{Source:} \texttt{threshold/<TARGET>/analyze\_calibration.py} already
produces \texttt{p\_bind\_distribution.png} and
\texttt{pIC50\_correlation.png}; combine into one two-panel figure with shared
styling rather than pasting two separate PNGs.
\textbf{Watch for:} save as PDF or 300\,dpi PNG. The default 150\,dpi in the
tutorial script is fine for a Colab preview and too low for print.
\end{figrec}

% \begin{figure}[t]
%   \centering
%   \includegraphics[width=\linewidth]{figures/fig2_calibration.pdf}
%   \caption{Boltz-2 calibration. \TODO{caption}}
%   \label{fig:calibration}
% \end{figure}

%-----------------------------------------------------------------------
\subsection{p38 MAPK14}
\label{sec:results:p38}

\begin{guide}[Per-iteration log blocks --- fill one in per iteration, the week it finishes]
Copy the block below for each new iteration. Keep them terse; they are notes to
yourself, and later the raw material for the prose. The "Observations" field is
the one that matters most and the one people skip --- write two honest sentences
about what the chemistry looked like and anything that surprised you.
\end{guide}

\subsubsection*{Iteration 1 --- unbiased prior}
\begin{guide}[Log]
\begin{tabular}{@{}p{0.28\linewidth}p{0.66\linewidth}@{}}
Date completed        & \FILL{} \\
Model $\Gmodel_1$     & \texttt{reinvent\_pubchem.prior} (no fine-tuning) \\
Generation job        & \FILL{SLURM job ID} \\
Docking job           & \FILL{SLURM job ID} \\
$|\Valid_1| / \Ngen$  & \FILL{} / 5{,}000 \\
$\tau$ used           & \FILL{} \\
$|\Passed_1|$, $\Kk$  & \FILL{} , \FILL{} \\
$\Tbar_1$ / $\Tstar_1$& \FILL{} / \FILL{} \\
Hits ($\ge 0.85$)     & \FILL{} \\
Scaffold count / $D_1$& \FILL{} / \FILL{} \\
Shared scaffolds with $\Actives$ & \FILL{} \\
Mean QED / SA         & \FILL{} / \FILL{} \\
Observations          & \FILL{Two sentences: what did the top-$\Kk$ chemistry
                        look like? Anything unexpected?} \\
Problems / deviations & \FILL{Timeouts, dropped molecules, anything you did
                        differently from the protocol.} \\
\end{tabular}
\smallskip

\emph{Expected at $t=1$:} zero or near-zero rediscovery hits and few or no
shared scaffolds. $\Gmodel_1$ has had no transfer learning and no knowledge of
either target. This row establishes the floor everything else is measured
against --- it is a result, not a failure.
\end{guide}

\subsubsection*{Iteration 2}
\begin{guide}[Log]
\begin{tabular}{@{}p{0.28\linewidth}p{0.66\linewidth}@{}}
Date completed        & \FILL{} \\
Model $\Gmodel_2$     & \FILL{checkpoint path from iteration 1 fine-tune} \\
Generation job        & \FILL{} \\
Docking job           & \FILL{} \\
$|\Valid_2| / \Ngen$  & \FILL{} / 5{,}000 \\
$|\Passed_2|$, $\Kk$  & \FILL{} , \FILL{} \\
$\Tbar_2$ / $\Tstar_2$& \FILL{} / \FILL{} \\
$\Delta \Tbar$ vs.\ $t=1$ & \FILL{} \\
Hits ($\ge 0.85$)     & \FILL{} \\
Scaffold count / $D_2$& \FILL{} / \FILL{} \\
New dominant scaffolds& \FILL{which ring systems appeared that weren't in
                        $\TopK_1$?} \\
Mean QED / SA         & \FILL{} / \FILL{} \\
Observations          & \FILL{} \\
Problems / deviations & \FILL{} \\
\end{tabular}
\end{guide}

\subsubsection*{Iteration 3}
\begin{guide}[Log --- copy this block for iterations 4, 5, \dots]
\begin{tabular}{@{}p{0.28\linewidth}p{0.66\linewidth}@{}}
Date completed        & \FILL{} \\
Model $\Gmodel_3$     & \FILL{} \\
$|\Valid_3| / \Ngen$  & \FILL{} / 5{,}000 \\
$|\Passed_3|$, $\Kk$  & \FILL{} , \FILL{} \\
$\Tbar_3$ / $\Tstar_3$& \FILL{} / \FILL{} \\
$\Delta \Tbar$ vs.\ $t=2$ & \FILL{} \\
Hits ($\ge 0.85$)     & \FILL{} \\
Scaffold count / $D_3$& \FILL{} / \FILL{} \\
Mean QED / SA         & \FILL{} / \FILL{} \\
Observations          & \FILL{} \\
Problems / deviations & \FILL{} \\
\end{tabular}
\end{guide}

\begin{yourturn}[Then write the p38 prose: 3--4 paragraphs]
Walk the reader down the columns of \cref{tab:master}, in this order:
\begin{enumerate}\itemsep0pt
  \item \textbf{Throughput and gating.} How many molecules survived validity
        filtering, how many cleared $\tau$, and how the pass rate moved across
        iterations. A rising pass rate is itself evidence the fine-tuning is
        working, independent of Tanimoto.
  \item \textbf{Score enrichment.} How the distribution of $S(m)$ shifted.
        Give median and IQR, not just the mean.
  \item \textbf{Convergence.} $\Tbar_t$ and $\Tstar_t$ against the $t=0$ null.
        State whether the trend is monotone and where it plateaus.
  \item \textbf{Hits.} How many, at which iteration, and what they resemble.
        If there are none, say that clearly and report the closest pairs
        instead --- that is still a quantitative result.
\end{enumerate}
\end{yourturn}

\TODO{p38 results prose}

\begin{figrec}[Figure 3 --- convergence curves (the paper's headline figure)]
$\Tbar_t$ (line, with bootstrap 95\% CI band) and $\Tstar_t$ (line, second
axis or second panel) versus iteration $t$, one colour per target on shared
axes, plus a flat dashed \emph{prior-baseline} line at the $t=0$ value and a
horizontal dashed line at $\thetahit = 0.85$.
\textbf{Why it matters:} this single panel is the paper's central claim. Make
it early, even with only two iterations of data --- looking at it will tell you
whether the loop is working long before the full run finishes.
\end{figrec}

\begin{figrec}[Figure 4 --- score distribution evolution]
Violin or ridgeline plot of $S(m)$ over $\TopK_t$, one violin per iteration,
faceted by target. Overlay the median. This is the plan's
"progressive enrichment" visualisation.
\end{figrec}

% \begin{figure}[t]
%   \centering
%   \includegraphics[width=0.85\linewidth]{figures/fig3_convergence.pdf}
%   \caption{Convergence of generated molecules toward the held-out actives.
%            \TODO{caption}}
%   \label{fig:convergence}
% \end{figure}

%-----------------------------------------------------------------------
\subsection{CDK2}
\label{sec:results:cdk2}

\begin{guide}[Same log structure --- copy the p38 blocks above]
Keep the fields identical to \cref{sec:results:p38} so the two targets stay
directly comparable. Resist the urge to track extra things for one target only;
asymmetric records make the cross-target comparison in
\cref{sec:results:cross} much harder to write.
\end{guide}

\subsubsection*{Iteration 1 --- unbiased prior}
\begin{guide}[Log]
\begin{tabular}{@{}p{0.28\linewidth}p{0.66\linewidth}@{}}
Date completed        & \FILL{} \\
Model $\Gmodel_1$     & \texttt{reinvent\_pubchem.prior} (no fine-tuning) \\
Generation job        & \FILL{} \\
Docking job           & \FILL{} \\
$|\Valid_1| / \Ngen$  & \FILL{} / 5{,}000 \\
$\tau$ used           & \FILL{} \\
$|\Passed_1|$, $\Kk$  & \FILL{} , \FILL{} \\
$\Tbar_1$ / $\Tstar_1$& \FILL{} / \FILL{} \\
Hits ($\ge 0.85$)     & \FILL{} \\
Scaffold count / $D_1$& \FILL{} / \FILL{} \\
Shared scaffolds with $\Actives$ & \FILL{} \\
Mean QED / SA         & \FILL{} / \FILL{} \\
Observations          & \FILL{} \\
Problems / deviations & \FILL{} \\
\end{tabular}
\end{guide}

\subsubsection*{Iteration 2}
\begin{guide}[Log]
\begin{tabular}{@{}p{0.28\linewidth}p{0.66\linewidth}@{}}
Date completed        & \FILL{} \\
Model $\Gmodel_2$     & \FILL{} \\
$|\Valid_2| / \Ngen$  & \FILL{} / 5{,}000 \\
$|\Passed_2|$, $\Kk$  & \FILL{} , \FILL{} \\
$\Tbar_2$ / $\Tstar_2$& \FILL{} / \FILL{} \\
$\Delta \Tbar$ vs.\ $t=1$ & \FILL{} \\
Hits ($\ge 0.85$)     & \FILL{} \\
Scaffold count / $D_2$& \FILL{} / \FILL{} \\
New dominant scaffolds& \FILL{} \\
Mean QED / SA         & \FILL{} / \FILL{} \\
Observations          & \FILL{} \\
Problems / deviations & \FILL{} \\
\end{tabular}
\end{guide}

\subsubsection*{Iteration 3}
\begin{guide}[Log --- copy for later iterations]
\begin{tabular}{@{}p{0.28\linewidth}p{0.66\linewidth}@{}}
Date completed        & \FILL{} \\
Model $\Gmodel_3$     & \FILL{} \\
$|\Valid_3| / \Ngen$  & \FILL{} / 5{,}000 \\
$|\Passed_3|$, $\Kk$  & \FILL{} , \FILL{} \\
$\Tbar_3$ / $\Tstar_3$& \FILL{} / \FILL{} \\
Hits ($\ge 0.85$)     & \FILL{} \\
Scaffold count / $D_3$& \FILL{} / \FILL{} \\
Mean QED / SA         & \FILL{} / \FILL{} \\
Observations          & \FILL{} \\
Problems / deviations & \FILL{} \\
\end{tabular}
\end{guide}

\begin{yourturn}[Then write the CDK2 prose: 3--4 paragraphs]
Same four-part structure as p38. Where CDK2 behaved differently, say so here in
factual terms and save the "why" for \cref{sec:discussion}. Note that CDK2's
benchmark set is roughly a third the size of p38's, which affects $\simmax$
directly: a smaller $\Actives$ means fewer chances for any generated molecule to
find a close neighbour. Mention this when comparing raw $\Tbar_t$ values across
targets --- otherwise the comparison is not apples to apples.
\end{yourturn}

\TODO{CDK2 results prose}

%-----------------------------------------------------------------------
\subsection{Chemical evolution across iterations}
\label{sec:results:chem}

\begin{yourturn}[Write 3--4 paragraphs. This is the interpretability contribution.]
Structure it as: scaffolds $\rightarrow$ functional groups $\rightarrow$
feasibility envelope.
\begin{itemize}\itemsep0pt
  \item \textbf{Scaffold trajectory.} Did $D_t$ (\cref{eq:scaffolddiv}) hold or
        collapse? Which scaffolds emerged, at which iteration, and do any
        resemble known pharmacophores for the target? Name the ring systems
        explicitly --- a reader who knows kinase chemistry will care whether an
        adenine mimetic or a hinge-binding motif appeared.
  \item \textbf{Functional groups.} Which $f_{g,t}$ trends were monotone?
        Interpret against binding-site character (a deep hydrophobic ATP site
        should not be enriching hydroxyls indefinitely). Only claim a trend if
        it holds over $\ge 3$ iterations.
  \item \textbf{Feasibility envelope.} Did MW, $c\log P$, QED and SA stay within
        bounds? Report the fraction of $\TopK_t$ violating each Lipinski
        threshold per iteration. If SA rose steadily, that is a real finding
        about over-optimisation and belongs in the abstract.
  \item \textbf{PAINS.} What fraction of top candidates carried a flag, and
        which catalogue entries.
\end{itemize}
\end{yourturn}

\TODO{chemical evolution prose}

\begin{figrec}[Figure 5 --- chemical trajectory (multi-panel, one row per metric)]
A grid: rows = \{scaffold diversity $D_t$, mean QED, mean SA score,
mean MW, mean $c\log P$\}, columns = \{p38, CDK2\}. Each panel is a line over
iteration with a shaded band for the known-actives reference range
(e.g.\ the actives' interquartile range as a horizontal band). The band is what
makes the figure readable: it turns "the number moved" into "the number moved
toward or away from real inhibitors".
\textbf{Source:} the per-iteration \texttt{*\_comparison\_profile.csv} files
from the chemical analysis notebooks.
\end{figrec}

\begin{figrec}[Figure 6 --- scaffold and functional-group evolution]
\textbf{(a)} Stacked bar of the top 8 Murcko scaffolds' share of $\TopK_t$ per
iteration, with structures drawn as insets for the 3 most frequent.
\textbf{(b)} Functional-group frequency $f_{g,t}$ trajectories, one line per
group, with the actives' frequency as a dashed horizontal reference per group.
\end{figrec}

\begin{figrec}[Figure 7 --- descriptor-space PCA]
Actives (grey, fitted space) with $\TopK_t$ overlaid per iteration in a
sequential colourmap (light $\rightarrow$ dark with increasing $t$). One panel
per target. If the loop is working, later iterations should visibly migrate
into the grey cloud. If they do not, that is also worth showing.
\textbf{Critical:} fit the PCA on actives only, transform both. Do not refit per
iteration or panels stop being comparable.
\end{figrec}

%-----------------------------------------------------------------------
\subsection{Rediscovery hit catalogue}
\label{sec:results:hits}

\begin{yourturn}[Write 1--2 paragraphs + fill \cref{tab:hits}]
For each molecule with $\simmax \ge \thetahit$: report the iteration it appeared,
its similarity, its nearest known active with that active's ChEMBL ID and
measured $\pICfifty$, and Boltz-2's predicted $\pICfifty$ for the generated
molecule. If there are no formal hits, retitle this subsection
"Closest generated--active pairs", report the top 5--10 by $\simmax$, and say
plainly that no molecule cleared $0.85$. That is an honest and perfectly
reportable outcome.
\end{yourturn}

\begin{table}[h]
\centering
\caption{Rediscovery hits and closest generated--active pairs.}
\label{tab:hits}
\small
\begin{tabular}{llcllcc}
\toprule
Target & Iter & $\simmax$ & Nearest active (ChEMBL ID) & Measured $\pICfifty$ &
Predicted $\pICfifty$ & PAINS \\
\midrule
\FILL{} & \FILL{} & \FILL{} & \FILL{} & \FILL{} & \FILL{} & \FILL{} \\
\FILL{} & \FILL{} & \FILL{} & \FILL{} & \FILL{} & \FILL{} & \FILL{} \\
\FILL{} & \FILL{} & \FILL{} & \FILL{} & \FILL{} & \FILL{} & \FILL{} \\
\bottomrule
\end{tabular}
\end{table}

\begin{figrec}[Figure 8 --- hit catalogue grid]
Side-by-side 2D structures: generated molecule (left) next to its nearest known
active (right), one row per hit, most similar first, annotated with $\simmax$
and both potency values. \textbf{Source:} \texttt{Draw.MolsToGridImage} in the
Part 2 analysis notebook already produces exactly this
(\texttt{*\_closest\_pairs.png}); increase \texttt{subImgSize} and save as SVG
or PDF for the paper.
\end{figrec}

%-----------------------------------------------------------------------
\subsection{Cross-target comparison}
\label{sec:results:cross}

\begin{yourturn}[Write 2 paragraphs]
This subsection is the generalisability claim, so it needs care.
Compare across targets: calibration quality (AUC, correlation), convergence rate
($\Delta \Tbar$ per iteration, not absolute $\Tbar$), hit iteration, and
scaffold behaviour. Explicitly control for the confounds:
different $|\Actives|$ (3{,}860 vs.\ 1{,}402), possibly different $\tau$, and
different pocket character. If the two targets behaved similarly, that supports
generalisation; if they diverged, identify which of these three confounds is the
likeliest explanation and say what experiment would separate them.
\end{yourturn}

\TODO{cross-target comparison}

%-----------------------------------------------------------------------
\subsection{Ablations}
\label{sec:results:ablation}

\begin{guide}[Optional --- include only what you actually ran]
Candidate ablations, roughly in order of value per unit of GPU time:
\begin{enumerate}\itemsep0pt
  \item \textbf{No-gate control:} rank by $-\yhat$ alone with no $\pbind$
        filter. Directly tests whether the gate matters. Cheapest and most
        informative.
  \item \textbf{Alternative scoring} (Options A and C from
        \cref{sec:methods:score}).
  \item \textbf{Gate sensitivity:} rerun with $\tau \pm 0.1$.
  \item \textbf{Corpus size:} $\Kk = 200$ vs.\ $500$ --- probes the
        exploitation/diversity trade-off.
  \item \textbf{Random-corpus control:} fine-tune on $\Kk$ \emph{randomly
        chosen} passing molecules instead of the top-$\Kk$. If $\Tbar_t$ rises
        anyway, your improvement is coming from fine-tuning on anything, not
        from the scoring signal. This is the single most convincing control in
        the whole paper --- run it if you run only one.
\end{enumerate}
Delete this subsection entirely if you ran none.
\end{guide}

\TODO{ablations, or delete this subsection}

%=======================================================================
\section{Discussion}
\label{sec:discussion}
%=======================================================================

\begin{guide}[Optional section --- keep if the venue expects it]
Four short blocks:
\begin{enumerate}\itemsep0pt
  \item \textbf{What the results mean.} One paragraph tying the convergence
        curves back to the claim in \cref{sec:intro:contrib}.
  \item \textbf{Limitations.} Predicted rather than measured affinity; possible
        overlap between Boltz-2's training data and the benchmark actives; two
        kinase targets do not test non-kinase generalisation; no synthesis or
        assay; single random seed if that is the case.
  \item \textbf{Comparison to prior work.} Where your numbers sit relative to
        the closed-loop papers reviewed in \cref{sec:intro:closedloop}.
  \item \textbf{Future work.} A non-kinase target; multi-objective scoring with
        SA/QED terms in the objective rather than only as monitoring; an
        experimental validation arm.
\end{enumerate}
Do not use the limitations paragraph to pre-empt criticism you could instead fix.
\end{guide}

\TODO{discussion, or delete}

%=======================================================================
\section*{Author contributions}
\begin{guide}[CRediT taxonomy --- fill in before submission]
Use the standard categories: Conceptualization, Methodology, Software,
Validation, Formal analysis, Investigation, Data curation, Writing --- original
draft, Writing --- review \& editing, Visualization, Supervision, Funding
acquisition. Every author needs at least one. Agree on this early; it is far
more awkward to negotiate the week before submission.
\end{guide}
\TODO{author contributions}

\section*{Acknowledgements}
\begin{guide}[Do not forget the ACCESS citation]
Computational work used the Anvil cluster at Purdue University through NSF
ACCESS allocation \texttt{bio220114}. ACCESS requires a specific
acknowledgement wording and citation --- look it up on the ACCESS site and use
their exact text.
\end{guide}
\TODO{acknowledgements + ACCESS allocation wording}

%=======================================================================
\bibliography{references}
%=======================================================================

\end{document}
